\documentclass[aspectratio=169,11pt]{beamer}
\usepackage{paperdeck}
\title[Subword Units (BPE)]{Neural Machine Translation of Rare Words with Subword Units}
\subtitle{Sennrich, Haddow and Birch (ACL 2016), arXiv:1508.07909}
\author[Paper summary]{Paper summary for the CS-781 capstone}
\date{}
\begin{document}

\begin{frame}\titlepage\end{frame}

\paperglance
{R.\ Sennrich, B.\ Haddow, A.\ Birch. \emph{Neural Machine Translation of Rare Words with Subword Units}. ACL 2016. University of Edinburgh.}
{Adapt the byte-pair-encoding compression algorithm to split words into frequent subword units, so a fixed-size vocabulary can represent \emph{any} word.}
{BPE is how our Phase-1 tokenizer is built (\code{scripts/02\_train\_tokenizer.py}). Every model in all seven phases reads BPE tokens.}
{Main deck: ``Step 02: a legal byte-level BPE tokenizer''. Glossary: ``Tokenization'' section.}

\begin{frame}{Roadmap}
  \begin{multicols}{2}\small\tableofcontents\end{multicols}
\end{frame}

% ------------------------------------------------------------------
\section{The problem}

\begin{frame}{Translation is an open-vocabulary problem}
  \begin{columns}[T]
    \begin{column}{0.52\textwidth}
      \small
      \begin{itemize}
        \item Neural models use a \textbf{fixed vocabulary}, typically 30,000--50,000 words.
        \item Real text keeps producing \textbf{new words}: names, compounds, inflections, loanwords.
        \item Words outside the vocabulary become \code{UNK}.
        \item Earlier fix: a \textbf{back-off dictionary}. Translate \code{UNK} by looking the source word up, or copy it.
      \end{itemize}
      \vspace{1mm}
      Problems with back-off:
      \begin{itemize}
        \item Words do not map one-to-one across languages.
        \item The model cannot \emph{generate} unseen words.
        \item Copying fails when alphabets differ, e.g.\ English to Russian.
      \end{itemize}
    \end{column}
    \begin{column}{0.44\textwidth}
      \begin{exampleblock}{The paper's example}
        \small
        German \emph{Abwasser$|$behandlungs$|$anlage} means ``sewage water treatment plant''.
        One long compound, three meaningful parts. A single fixed-length vector for the whole word is a poor fit; its parts are what translate.
      \end{exampleblock}
      \begin{exampleblock}{A legal example}
        \small
        \emph{Vakalatnama}, \emph{Panchnama}, \emph{Mukhtarnama} share the suffix \emph{-nama} (``document''). A word-level vocabulary may miss the rarer ones entirely.
      \end{exampleblock}
    \end{column}
  \end{columns}
\end{frame}

\begin{frame}{Key intuition: rare words are made of translatable pieces}
  \begin{columns}[T]
    \begin{column}{0.52\textwidth}
      \small
      Many word types can be translated by a competent translator even if new, from known smaller units:
      \begin{itemize}
        \item \textbf{Names}: copy or transliterate. \emph{Barack Obama} becomes Russian \emph{Barak Obama} (in Cyrillic).
        \item \textbf{Cognates and loanwords}: regular character changes. \emph{claustrophobia}, \emph{Klaustrophobie}.
        \item \textbf{Morphologically complex words}: translate the parts. \emph{solar system}, \emph{Sonnen$|$system}.
      \end{itemize}
    \end{column}
    \begin{column}{0.44\textwidth}
      \begin{block}{Evidence from the paper}
        \small
        100 rare German tokens (outside the top 50,000) were analysed:
        \vspace{1mm}

        \begin{tabular}{@{}lr@{}}
          \toprule
          Category & Count\\
          \midrule
          Compounds & 56\\
          Names & 21\\
          Loanwords, common origin & 6\\
          Transparent affixation & 5\\
          Number / code identifier & 2\\
          \bottomrule
        \end{tabular}
        \vspace{1mm}

        Most rare words are translatable from their parts.
      \end{block}
    \end{column}
  \end{columns}
\end{frame}

\begin{frame}{The trade-off: vocabulary size vs sequence length}
  \begin{columns}[T]
    \begin{column}{0.46\textwidth}
      \small
      \begin{itemize}
        \item \textbf{Words}: short sequences, but a huge vocabulary and unknown words.
        \item \textbf{Characters}: tiny vocabulary and no unknowns, but very long sequences. Attention must span far more positions.
        \item \textbf{Goal}: a compact vocabulary \emph{and} short text, with no unknown words.
      \end{itemize}
      \vspace{1mm}
      Earlier segmenters from statistical MT, such as compound splitting and Morfessor, only moderately shrink the vocabulary and still leave unknown words.
    \end{column}
    \begin{column}{0.5\textwidth}
      \footnotesize
      German training data (about 100M tokens):
      \vspace{1mm}

      \begin{tabular}{@{}lrrr@{}}
        \toprule
        Segmentation & Tokens & Types & UNK\\
        \midrule
        none (words) & 100 m & 1,750,000 & 1079\\
        characters & 550 m & 3,000 & 0\\
        char.\ bigrams & 306 m & 20,000 & 34\\
        compound splitting & 102 m & 1,100,000 & 643\\
        Morfessor & 109 m & 544,000 & 237\\
        \textbf{BPE} & \textbf{112 m} & \textbf{63,000} & \textbf{0}\\
        BPE (joint) & 111 m & 82,000 & 32\\
        \bottomrule
      \end{tabular}
      \vspace{1mm}

      {\scriptsize UNK = unknown tokens in newstest2013.}
      \vspace{1mm}

      BPE makes text only 12\% longer than words, with a vocabulary 28$\times$ smaller and \textbf{no unknowns}.
    \end{column}
  \end{columns}
\end{frame}

% ------------------------------------------------------------------
\section{The basic idea: byte-pair encoding}

\begin{frame}{From compression to segmentation}
  \begin{columns}[T]
    \begin{column}{0.52\textwidth}
      \small
      \textbf{Original BPE} (Gage, 1994) is a data-compression method: repeatedly replace the most frequent pair of bytes with a new, unused byte.
      \vspace{1mm}

      \textbf{Adaptation for words}:
      \begin{enumerate}
        \item Write every word as characters plus an end-of-word marker \code{</w>}.
        \item Count all adjacent symbol pairs, weighted by word frequency.
        \item Merge the most frequent pair (A, B) into a new symbol AB everywhere.
        \item Repeat for a chosen number of merges.
      \end{enumerate}
      Vocabulary = characters + merges; the \textbf{number of merges is the only hyperparameter}.
    \end{column}
    \begin{column}{0.44\textwidth}
      \begin{block}{Design details}
        \small
        \begin{itemize}
          \item Pairs never cross word boundaries, so BPE can run on a word-frequency dictionary, which is fast.
          \item The \code{</w>} marker lets the original text be restored exactly.
          \item Frequent words end up as one symbol; rare words as several pieces.
          \item Unlike Huffman codes, the pieces are \emph{readable} subwords, so the network can reuse them.
        \end{itemize}
      \end{block}
    \end{column}
  \end{columns}
\end{frame}

\begin{frame}{Easy example: the paper's toy vocabulary}
  \begin{columns}[T]
    \begin{column}{0.48\textwidth}
      \small
      Word counts:
      \code{low}$\times$5, \code{lower}$\times$2, \code{newest}$\times$6, \code{widest}$\times$3.
      \vspace{1mm}

      \begin{tabular}{@{}clr@{}}
        \toprule
        Step & Merge & Count\\
        \midrule
        1 & e + s $\to$ es & 9\\
        2 & es + t $\to$ est & 9\\
        3 & est + \code{</w>} $\to$ est\code{</w>} & 9\\
        4 & l + o $\to$ lo & 7\\
        5 & lo + w $\to$ low & 7\\
        6 & n + e $\to$ ne & 6\\
        7 & ne + w $\to$ new & 6\\
        8 & new + est\code{</w>} $\to$ newest\code{</w>} & 6\\
        \bottomrule
      \end{tabular}
      \vspace{1mm}

      {\scriptsize Counts: ``e s'' appears in newest (6) and widest (3), so 9.}
    \end{column}
    \begin{column}{0.48\textwidth}
      \begin{exampleblock}{Applying the merges to text}
        \small
        After 10 merges the vocabulary gives:
        \begin{itemize}
          \item \code{newest} $\to$ \code{newest</w>}: one symbol.
          \item \code{widest} $\to$ \code{wi d est</w>}.
          \item \code{lower} $\to$ \code{low e r </w>}.
        \end{itemize}
        A \textbf{new} word \code{lowest}, never seen, becomes \code{low est</w>}, two known pieces.
      \end{exampleblock}
      \begin{exampleblock}{Legal flavour}
        \small
        If ``appeal'', ``appellant'' and ``appellate'' are frequent, BPE learns \code{appe}, \code{appell}, and so on. The rare \emph{appellee} then splits into known pieces like \code{appell} + \code{ee}.
      \end{exampleblock}
    \end{column}
  \end{columns}
\end{frame}

\begin{frame}[fragile]{The algorithm in a few lines}
  \begin{columns}[T]
    \begin{column}{0.55\textwidth}
      \scriptsize
\begin{verbatim}
def get_stats(vocab):             # count adjacent pairs
    pairs = defaultdict(int)
    for word, freq in vocab.items():
        s = word.split()
        for i in range(len(s) - 1):
            pairs[s[i], s[i+1]] += freq
    return pairs

vocab = {'l o w </w>': 5, 'l o w e r </w>': 2,
         'n e w e s t </w>': 6, 'w i d e s t </w>': 3}
for i in range(num_merges):
    pairs = get_stats(vocab)
    best = max(pairs, key=pairs.get)  # most frequent
    vocab = merge_vocab(best, vocab)  # replace everywhere
\end{verbatim}
      {\scriptsize Adapted from the paper's Algorithm 1.}
    \end{column}
    \begin{column}{0.41\textwidth}
      \begin{block}{At test time}
        \small
        Split a word into characters, then apply the learned merges \emph{in order}.
        This works for \emph{any} word. The only possible unknowns are unseen characters.
      \end{block}
      \begin{block}{Cost}
        \small
        Learning is linear in dictionary size per merge. Real implementations index pairs and update counts incrementally.
      \end{block}
    \end{column}
  \end{columns}
\end{frame}

\begin{frame}{Joint BPE: consistent pieces in both languages}
  \begin{columns}[T]
    \begin{column}{0.5\textwidth}
      \small
      Two options:
      \begin{itemize}
        \item \textbf{Separate BPE}: learn one encoding per language. More compact, and every piece is seen in its own language.
        \item \textbf{Joint BPE}: learn one encoding on the union of both vocabularies. Names are split \emph{the same way} on both sides, so they are easier to copy or transliterate.
      \end{itemize}
      For English--Russian, Russian is transliterated to Latin script for learning, and the merges are mapped back to Cyrillic.
    \end{column}
    \begin{column}{0.46\textwidth}
      \begin{exampleblock}{Why consistency matters (paper's example)}
        \small
        Name \emph{Mirzayeva}; Russian shown transliterated:
        \begin{itemize}
          \item Separate BPE: \code{Mirz|ayeva} $\to$ \code{Mir|za|eva}. The pieces do not line up.
          \item Joint BPE: \code{Mir|za|yeva} $\to$ \code{Mir|za|eva}. Piece-by-piece transliteration.
        \end{itemize}
        Inconsistent splits in training data taught the separate system errors like \emph{rak} $\to$ \emph{pra}.
      \end{exampleblock}
    \end{column}
  \end{columns}
\end{frame}

% ------------------------------------------------------------------
\section{Experiments}

\begin{frame}{Setup}
  \begin{columns}[T]
    \begin{column}{0.5\textwidth}
      \small
      \begin{itemize}
        \item \textbf{Task}: WMT 2015, English$\to$German (4.2M sentence pairs, about 100M tokens) and English$\to$Russian (2.6M pairs, about 50M tokens).
        \item \textbf{Model}: attention-based RNN encoder--decoder (Bahdanau et al., 2015). Hidden size 1000, embeddings 620.
        \item \textbf{Training}: Adadelta, batch 80, about 7 days per model; beam size 12.
        \item \textbf{Metrics}: BLEU, chrF3 (character n-gram F-score), and unigram F1 split into all, \textbf{rare} and \textbf{OOV} words.
      \end{itemize}
    \end{column}
    \begin{column}{0.46\textwidth}
      \begin{block}{Systems compared}
        \small
        \begin{itemize}
          \item \textbf{WUnk}: words; unknowns stay \code{UNK}.
          \item \textbf{WDict}: words plus back-off dictionary (baseline).
          \item \textbf{C2-50k}: 50k frequent words, character bigrams for the rest.
          \item \textbf{BPE-60k}: separate BPE.
          \item \textbf{BPE-J90k}: joint BPE.
        \end{itemize}
        No subword system uses a back-off dictionary.
      \end{block}
    \end{column}
  \end{columns}
\end{frame}

\begin{frame}{Results: English$\to$German (newstest2015)}
  \footnotesize
  \begin{tabular}{@{}lcccccc@{}}
    \toprule
    & \multicolumn{2}{c}{BLEU} & \multicolumn{2}{c}{chrF3} & \multicolumn{2}{c}{unigram F1 (\%), ensemble}\\
    System & single & ens-8 & single & ens-8 & rare & OOV\\
    \midrule
    Syntax-based SMT & 24.4 & -- & 55.3 & -- & -- & --\\
    WUnk & 20.6 & 22.8 & 47.2 & 48.9 & 20.4 & 0.0\\
    WDict (baseline) & 22.0 & 24.2 & 50.5 & 52.4 & 36.8 & 36.8\\
    C2-50k & \textbf{22.8} & \textbf{25.3} & 51.9 & 53.5 & 40.5 & 30.9\\
    BPE-60k & 21.5 & 24.5 & 52.0 & 53.9 & 40.9 & 29.3\\
    BPE-J90k & \textbf{22.8} & 24.7 & 51.7 & \textbf{54.1} & \textbf{41.8} & 33.6\\
    \bottomrule
  \end{tabular}
  \vspace{2mm}
  \small
  \begin{itemize}
    \item Rare-word unigram F1 rises from 36.8 to \textbf{41.8}.
    \item For English$\to$German, copying unknown names works well, so the dictionary's OOV score stays high.
    \item Subword ensembles beat the baseline by 0.3--1.1 BLEU on this pair.
  \end{itemize}
\end{frame}

\begin{frame}{Results: English$\to$Russian (different alphabets)}
  \footnotesize
  \begin{tabular}{@{}lcccccc@{}}
    \toprule
    & \multicolumn{2}{c}{BLEU} & \multicolumn{2}{c}{chrF3} & \multicolumn{2}{c}{unigram F1 (\%), ensemble}\\
    System & single & ens-8 & single & ens-8 & rare & OOV\\
    \midrule
    Phrase-based SMT & 24.3 & -- & 53.8 & -- & -- & --\\
    WUnk & 18.8 & 22.4 & 46.5 & 49.9 & 25.2 & 0.0\\
    WDict (baseline) & 19.1 & 22.8 & 47.5 & 51.0 & 26.5 & 6.6\\
    C2-50k & \textbf{20.9} & \textbf{24.1} & 49.0 & 51.6 & 27.8 & 17.4\\
    BPE-60k & 20.5 & 23.6 & \textbf{49.8} & 52.7 & 29.7 & 15.6\\
    BPE-J90k & 20.4 & \textbf{24.1} & 49.7 & \textbf{53.0} & 29.7 & \textbf{18.3}\\
    \bottomrule
  \end{tabular}
  \vspace{2mm}
  \small
  \begin{itemize}
    \item Copying fails across alphabets, so the dictionary gets only 6.6 OOV F1. Subwords reach \textbf{18.3}, because they learn to \emph{transliterate}.
    \item Joint BPE gains 1.3 BLEU and 2.0 chrF3 over the baseline.
  \end{itemize}
\end{frame}

\begin{frame}{Analysis: rare words and what the models produce}
  \begin{columns}[T]
    \begin{column}{0.5\textwidth}
      \small
      \textbf{By word frequency}: all systems are similar on the 50,000 most frequent words.
      Between ranks 50,000 and 500,000, representing words as \emph{subwords} beats representing them as rare whole words.
      Rank 50,000 means about 60 training occurrences; rank 500,000 means about 2.
      \vspace{1mm}

      \textbf{Lesson}: subword units are \emph{less sparse} than rare words. A \emph{smaller} vocabulary with more subword splitting can \emph{improve} accuracy.
    \end{column}
    \begin{column}{0.46\textwidth}
      \begin{exampleblock}{Translations from the paper}
        \scriptsize
        \emph{health research institutes} $\to$ reference \emph{Gesundheitsforschungsinstitute}
        \begin{itemize}
          \item WDict: \emph{Forschungsinstitute} (drops ``health'').
          \item BPE-J90k: \emph{Gesundheits$|$forsch$|$ungsin$|$stitute}. Correct, even though the splits ignore morpheme boundaries.
        \end{itemize}
        \emph{rakfisk} $\to$ reference \emph{rakfiska} (Cyrillic)
        \begin{itemize}
          \item WDict: copies \emph{rakfisk} unchanged.
          \item BPE-J90k: \emph{rak$|$f$|$iska}. Transliterated correctly.
        \end{itemize}
      \end{exampleblock}
    \end{column}
  \end{columns}
\end{frame}

% ------------------------------------------------------------------
\section{Assessment and outlook}

\begin{frame}{Connection to our capstone}
  \begin{columns}[T]
    \begin{column}{0.5\textwidth}
      \small
      \begin{itemize}
        \item Our tokenizer is \textbf{byte-level BPE} (the GPT-2 variant): it starts from 256 bytes instead of characters, so even unseen characters cannot be unknown.
        \item \code{tokenizer.vocab\_size} (default 32,000) sets the number of merges, the paper's only hyperparameter.
        \item Our \textbf{fertility} metric (tokens per word) is the paper's ``text size'' axis.
      \end{itemize}
    \end{column}
    \begin{column}{0.46\textwidth}
      \begin{exampleblock}{Experiments the paper suggests}
        \small
        \begin{itemize}
          \item Vocabulary size sweep: 16k, 32k, 48k. The paper found a smaller vocabulary can help rare words.
          \item Legal BPE vs GPT-2's web BPE: which splits legal words better?
          \item Check rare legal terms and Indian names: how many pieces do they take?
          \item For NER (Phase 2): do BPE splits break entity boundaries?
        \end{itemize}
      \end{exampleblock}
    \end{column}
  \end{columns}
\end{frame}

\begin{frame}{Strengths and limitations}
  \begin{columns}[T]
    \begin{column}{0.48\textwidth}
      \begin{exampleblock}{Strengths}
        \small
        \begin{itemize}
          \item Simple, fast, language-independent.
          \item True open vocabulary with a compact, fixed symbol set.
          \item One hyperparameter.
          \item Removes the need for back-off dictionaries.
          \item Became the default tokenizer family for modern language models.
        \end{itemize}
      \end{exampleblock}
    \end{column}
    \begin{column}{0.48\textwidth}
      \begin{alertblock}{Limitations}
        \small
        \begin{itemize}
          \item Greedy and frequency-based: splits often ignore morphology.
          \item Vocabulary size was chosen ``somewhat arbitrarily'', by the authors' own account.
          \item Deterministic segmentation: a word is always split the same way, so the model never sees alternatives.
          \item Separate BPE can split names inconsistently across languages.
          \item BLEU and chrF3 disagreed on the best system; gains on rare words are hard to see in aggregate metrics.
        \end{itemize}
      \end{alertblock}
    \end{column}
  \end{columns}
\end{frame}

\begin{frame}{Future research directions (1): from the paper and after}
  \small
  \begin{tabularx}{\textwidth}{@{}lX@{}}
    \toprule
    Direction & Details\\
    \midrule
    Learn the vocabulary size automatically & Proposed in the paper: the best size should depend on the language pair and data size\\
    Bilingually informed segmentation & Proposed in the paper: create more alignable units, while working without the target text at runtime\\
    Byte-level BPE & GPT-2 (Radford et al., 2019) runs BPE on bytes, so no symbol can ever be unknown\\
    Probabilistic segmentation & Unigram LM tokenizer and SentencePiece (Kudo, 2018; Kudo and Richardson, 2018)\\
    Segmentation as regularisation & BPE-dropout (Provilkov et al., 2020) randomly skips merges during training\\
    Tokenizer-free models & ByT5 (Xue et al., 2022) and CANINE (Clark et al., 2022) read bytes or characters directly\\
    \bottomrule
  \end{tabularx}
\end{frame}

\begin{frame}{Future research directions (2): open questions}
  \begin{columns}[T]
    \begin{column}{0.48\textwidth}
      \begin{block}{Open in general}
        \small
        \begin{itemize}
          \item What makes a tokenizer ``good'' for a model, beyond compression?
          \item Fairness across languages: the same content can cost several times more tokens in some languages than in English (Petrov et al., 2023).
          \item Handling numbers and dates consistently.
          \item Reducing variance between training runs, which the paper flags as an open problem for NMT.
        \end{itemize}
      \end{block}
    \end{column}
    \begin{column}{0.48\textwidth}
      \begin{exampleblock}{Open for legal NLP}
        \small
        \begin{itemize}
          \item Keeping citations such as ``(2019) 5 SCC 123'' and section numbers such as ``302'' intact or predictably split.
          \item Indian names and transliterated terms in English judgments.
          \item Domain-adaptive vocabularies: add legal merges to a general tokenizer.
          \item Character offsets for NER: subword boundaries must map back to exact characters.
        \end{itemize}
      \end{exampleblock}
    \end{column}
  \end{columns}
\end{frame}

\begin{frame}{Summary}
  \begin{itemize}
    \item Word-level NMT cannot handle an open vocabulary; back-off dictionaries are a brittle patch.
    \item \textbf{BPE}: start from characters and repeatedly merge the most frequent adjacent pair. Frequent words become single units; rare words become known pieces.
    \item On German data: 63k types instead of 1.75M, text only 12\% longer, and no unknown tokens.
    \item Better rare-word translation (EN$\to$DE rare unigram F1 36.8 to 41.8), transliteration across alphabets, and up to 1.1 (EN$\to$DE) and 1.3 (EN$\to$RU) BLEU.
    \item Joint BPE keeps segmentations consistent across languages.
    \item For us: the foundation of the Phase-1 tokenizer, and a source of tokenizer experiments.
  \end{itemize}
\end{frame}

\begin{frame}{Reference}
  \small
  Sennrich, R., Haddow, B., and Birch, A. (2016). Neural Machine Translation of Rare Words with Subword Units. In \emph{Proceedings of ACL 2016}. arXiv:1508.07909.
  \vspace{3mm}

  {\footnotesize Other work cited: Gage (1994) A New Algorithm for Data Compression; Bahdanau et al.\ (2015); Radford et al.\ (2019) GPT-2; Kudo (2018) Subword Regularization; Kudo and Richardson (2018) SentencePiece; Provilkov et al.\ (2020) BPE-Dropout; Xue et al.\ (2022) ByT5; Clark et al.\ (2022) CANINE; Petrov et al.\ (2023) Language Model Tokenizers Introduce Unfairness Between Languages.}
\end{frame}

\end{document}
